\section{A mechanical construction with a finite auxiliary history}
\label{sec:mechanical}

The construction uses an indefinite separable Hamiltonian.  Two new
families of compact ovals occur near the coordinate axes after each
quadratic fold.  The periods of each source form vanish on the opposite family by
reflection.  The additional logarithm comes from an auxiliary mechanical
history: at a given global stage, all earlier auxiliary birth levels
contribute at the same leading order after their real copies are counted.
For arbitrary degrees, the potentials are padded by small positive even
monomials.  We fix the complete finite list of future auxiliary budgets
before making these perturbations.

The period calculation uses the monomial powers of
\cite[p.~445]{LiLiLlibreZhang2001} and the outer-period asymptotics of
\cite[Lemma~2.3]{HanRomanovski2011}. We give the required version for
arbitrary even potentials, since each fold changes the potential.

\begin{lemma}[Mechanical outer periods]\label{lem:mechanical-rank}
Let $W$ be an even real polynomial of degree $2s$ with positive leading
coefficient.  On the outer period annulus of
$F(z,w)=z^2/2+W(w)$, the periods
\[
 I_{p,q}(h)=\oint_{F=h}z^{2p+1}w^{2q}\,dw,
 \qquad p\geq0,\quad 0\leq q<s,
\]
are linearly independent for every finite collection of distinct pairs
$(p,q)$.  A collection of $M$ such forms therefore has a linear
combination with $M-1$ simple period zeros in any prescribed nonempty
open interval of the outer annulus.
\end{lemma}

\begin{proof}
Write the leading term of $W$ as $c w^{2s}$, $c>0$.  Scaling $w$ by
$(h/c)^{1/(2s)}$ in the integral between the two outer turning points gives
\[
 I_{p,q}(h)\sim C_{p,q}
 h^{p+1/2+(2q+1)/(2s)},\qquad C_{p,q}\ne0.
\]
For a consistent oval orientation the constants have a fixed sign:
apart from that sign, the scaled integral is a positive multiple of
$\int_{-1}^{1}(1-u^{2s})^{p+1/2}u^{2q}\,du$.
The exponents are distinct, since their equality would imply
$p-p'=-(q-q')/s$ with $|q-q'|<s$.  The term with greatest exponent
excludes every nontrivial finite linear relation.  The periods are
analytic on the connected outer energy interval, so their independence
holds on each open subinterval.  Lemma~\ref{lem:jets} gives the stated
simple zeros.  Each zero can be assigned a compact regular annular
neighborhood and a positive $C^1$ persistence margin.
\end{proof}

Reflection cancellation of inherited periods is already used in the
classical fourfold construction \cite[equation~(3.3)]{ChristopherLloyd1995}.
The following observation records its persistence under the small
Hamiltonian changes needed here.

\begin{lemma}[Symmetry on continued ovals]\label{lem:continued-symmetry}
Let $F$ be a smooth planar Hamiltonian satisfying $F\circ R=F$ on a
neighborhood of a compact regular annular family, where $R$ is a
coordinate reflection preserving each oval of the family.
If $G\circ R=G$ and $G$ is sufficiently close to $F$ in $C^3$ on a
slightly larger annulus, the nearby regular ovals of $G$ are also
preserved by $R$.  If $R^*\alpha=\alpha$, the period of $\alpha$ on
this entire nearby family is identically zero.
\end{lemma}

\begin{proof}
Use the flow of $\nabla F/|\nabla F|^2$ to write an invariant annular
collar as $(\theta,h)$ with $F=h$.  The flow is equivariant under $R$.
For $G$ close to $F$, the equation $G(\theta,u)=h$ has a unique solution
$u=u(\theta,h)$ in this collar, by the implicit function theorem.
Uniqueness and $G\circ R=G$ make its closed normal graph invariant under
$R$.  It is a regular closed level component, hence an actual periodic
oval.  A reflection reverses its oval orientation, and consequently
\[
 \oint\alpha=\oint R^*\alpha=-\oint\alpha=0.
\]
This holds for every energy in an interval, so it includes the energy
derivative of the period.  The same collar construction gives the
$C^1$ continuity of periods used below, as in
Lemma~\ref{lem:persistence}.
\end{proof}

\begin{proposition}[Mechanical renewal]\label{prop:mechanical-certificate}
Fix an integer $n\geq2$ and put
\[
 \begin{gathered}
 r=\lfloor\log_2 n\rfloor,\qquad
 n_i=\left\lfloor\frac{n}{2^{r-i}}\right\rfloor
 \quad(0\leq i\leq r),\\
 \eta_i=n_i-2n_{i-1}\in\{0,1\}\quad(1\le i\le r).
 \end{gathered}
\]
For $1\leq i\leq r$ and $0\leq k<i$, set
$e_{i,k}=n_{i-1}\bmod 2^k$ and
\begin{equation}\label{eq:mechanical-birth}
 K_i=i(n_i^2-\eta_i)
       -4\sum_{k=0}^{i-1}(e_{i,k}^2+\eta_i e_{i,k})
       -4(2^i-1).
\end{equation}
There are even polynomials $U_r,V_r$ of degree $2n$ with positive
leading coefficients and a polynomial one-form $\omega_r$ of degree
at most $2n-1$ such that its period functions on compact regular annuli
of $H_r(x,y)=U_r(x)-V_r(y)$ have at least
\begin{equation}\label{eq:mechanical-total}
 L(n)=\sum_{i=1}^{r}4^{r-i}K_i
\end{equation}
simple zeros with pairwise disjoint annular neighborhoods.
In particular, there exists a real polynomial vector field of degree at
most $2n-1$ having at least $L(n)$ hyperbolic limit cycles. The bound satisfies
\begin{equation}\label{eq:finitebound}
 L(n)\ge\frac12n^2r^2+\frac16n^2r-\frac{22}{3}n^2.
\end{equation}
\end{proposition}

\begin{proof}
We construct the Hamiltonian and form by a finite induction, retaining
the auxiliary forms needed at later stages. The estimate
\eqref{eq:finitebound} is proved in Section~\ref{sec:counting}.

\paragraph{Auxiliary budgets.}
For each future global stage $t\in\{1,\ldots,r\}$, we build an auxiliary
form through levels $j=0,\ldots,t-1$ for each of the potentials $U_j$
and $V_j$.  Denote either such form by
\[
 \Omega_j^{(t)}=\sum_{p\geq0}z^{2p+1}R_p(w)\,dw.
\]
The final pullback into a stage-$t$ axis source has degree at most
$2(p+\deg R_p+1)$, so its endpoint budget is
$p+\deg R_p\leq n_t-2$.  Reversing the remaining folds gives the
budget at level $j$:
\begin{equation}\label{eq:mechanical-future-budget}
 p+2^{t-1-j}(\deg R_p+1)-1\leq n_t-2
 \quad\text{for each nonzero }R_p.
\end{equation}
A fold $w\mapsto w^2-b$ replaces $R_p$ by
$2wR_p(w^2-b)$, with
$\deg(2wR_p(w^2-b))+1=2(\deg R_p+1)$ for $R_p\ne0$.
Thus it preserves~\eqref{eq:mechanical-future-budget} at
the next level.  A birth at level $\ell$ uses exactly the monomials
\begin{equation}\label{eq:mechanical-birth-modes}
 z^{2p+1}w^{2q}\,dw,\qquad
 0\leq q<n_\ell,\qquad
 0\leq p\leq n_t-1-2^{t-1-\ell}(2q+1).
\end{equation}
They satisfy the budget.  Their periods are independent by
Lemma~\ref{lem:mechanical-rank}.

Set $U_0(w)=V_0(w)=w^2$ and $\omega_0=0$.  For every label $t$,
choose the level-zero auxiliary form in the span of the monomials
\eqref{eq:mechanical-birth-modes}, with one fewer selected simple period
zero than the dimension of that span, and with compact regular annular
neighborhoods around those zeros.  For a one-dimensional
space this requirement is empty.  The complete list of labels is fixed
at this point.

Suppose that stage $i-1$ has been completed.  Its data are the actual
potentials, one global form and all its selected zeros, and the
auxiliary forms $\Omega_{i-1}^{(t)}$ for $t\geq i$, separately for $U$
and $V$.  Every selected zero has a compact regular annular neighborhood
and a positive $C^1$ margin.  The list is finite.

\paragraph{Real copies and the two mechanical sources.}
Write $U=U_{i-1}$, $V=V_{i-1}$, and first consider
\[
 \widehat H_i(x,y)=U(x^2-a)-V(y^2-b),\qquad
 \Phi(x,y)=(x^2-a,y^2-b).
\]
For sufficiently large positive $a,b$, the four real inverse branches
of $\Phi$ cover all selected old global annuli.  They give four
diffeomorphic copies of each old simple period zero for
$\Phi^*\omega_{i-1}$, up to orientation.  The maps $(z,w)\mapsto(z,w^2-a)$ and
$(z,w)\mapsto(z,w^2-b)$ likewise have two diffeomorphic inverse branches
on the selected $U$- and $V$-auxiliary annuli, respectively. Pulling back
the auxiliary forms preserves their period functions up to orientation,
so every selected auxiliary zero has two simple copies.

The current vertical source uses the already fixed $V$-auxiliary form
$\Omega_{i-1,V}^{(i)}$.  Put
\[
 c_a=\sqrt{2U'(a)},\qquad z=c_a x,\qquad w=y^2-b.
\]
On either real $y$ branch the Hamiltonian deficit is exactly
\begin{equation}\label{eq:mechanical-limit}
 U(a)-\widehat H_i=K_a(z)+V(w),\qquad
 K_a(z)=U(a)-U\left(-a+\frac{z^2}{2U'(a)}\right).
\end{equation}
Since $U$ is even with positive leading coefficient, $U'(a)>0$ for
large $a$.  Taylor expansion gives $K_a(z)=z^2/2+o(1)$ in every fixed
$C^j$ norm on compact $z$ intervals.  More precisely, if $d=\deg U$,
the coefficient of $z^{2k}$ for $k\geq2$ is
$O(a^{-(k-1)d})$. Every regular mechanical oval is individually fixed
by $z\mapsto-z$: over its $w$ interval it consists of the two graphs
$z=\pm\sqrt{2(h-V(w))}$ joined at their turning points. This holds
also for a potential well centered at a nonzero value of $w$. The
$C^j$ convergence of $K_a$ to $z^2/2$ gives continuations of the auxiliary
ovals and their simple period zeros to \eqref{eq:mechanical-limit}, by
Lemma~\ref{lem:persistence}. Lemma~\ref{lem:continued-symmetry} preserves
the reflection symmetry of these continued ovals.  Pull back the fixed auxiliary form to
the polynomial source
\begin{equation}\label{eq:mechanical-vertical}
 \alpha_V=
 \sum_p(c_a x)^{2p+1}R_p(y^2-b)\,2y\,dy.
\end{equation}
Its periods on the two branches are $C^1$ close, up to orientation, to
the selected auxiliary periods.  This produces two copies of every
auxiliary simple zero.

Interchanging $U,x,a$ with $V,y,b$ gives a horizontal source
\begin{equation}\label{eq:mechanical-horizontal}
 \alpha_U=
 \sum_p(c_b y)^{2p+1}\widetilde R_p(x^2-a)\,2x\,dx,
 \qquad c_b=\sqrt{2V'(b)},
\end{equation}
with two copies of the corresponding $U$-auxiliary zeros.
The vertical approximation uses only large $a$ and fixed previous
$V$-annuli; $b$ only has to cover their $w$ ranges.  The horizontal
approximation uses large $b$ and fixed previous $U$-annuli, while $a$
covers their $w$ ranges.  Thus these are compatible lower bounds on
$a,b$.  In particular, no approximation is required near the new outer
barriers created by the present folds.

All these annuli can be made disjoint.  Old global copies have
$|x|\sim\sqrt a$, $|y|\sim\sqrt b$; the vertical sources have
$x=O(1/c_a)$ and $|y|\sim\sqrt b$; the horizontal sources have
$y=O(1/c_b)$ and $|x|\sim\sqrt a$.  Within a source, disjointness
comes from the selected auxiliary annular neighborhoods.  The ovals
are genuine closed level components, obtained in regular annular
collars; properness of the indefinite Hamiltonian is not needed.

\paragraph{Positive even padding.}
Fix such $a,b$, and define
\begin{equation}\label{eq:mechanical-padded-potentials}
 U_i(x)=U(x^2-a)+\epsilon x^{2n_i},\qquad
 V_i(y)=V(y^2-b)+\delta y^{2n_i},\qquad H_i=U_i-V_i.
\end{equation}
If $\eta_i=0$, set $\epsilon=\delta=0$.  If $\eta_i=1$, choose both
coefficients strictly positive and sufficiently small.  The resulting
potentials have exact degree $2n_i$, are even, and have positive leading
coefficients.

The finite conditions on this choice are the persistence of the four
old global copies, the two auxiliary copies for every label $t>i$,
and the current vertical and horizontal source zeros for the fixed
forms~\eqref{eq:mechanical-vertical} and
\eqref{eq:mechanical-horizontal}.  All are conditions on already fixed
compact sets and forms.  In vertical source coordinates the new error
is exactly
\[
 -\epsilon(z/c_a)^{2n_i}+\delta(w+b)^{n_i}.
\]
For fixed $a,b,c_a$ its $C^3$ norm on each selected compact rectangle
tends to zero with $(\epsilon,\delta)$.  In horizontal coordinates the
error is $\epsilon(w+a)^{n_i}-\delta(z/c_b)^{2n_i}$ and has the same
property.  The auxiliary and old global errors also tend to zero on
their fixed collars.  Lemma~\ref{lem:persistence} therefore imposes
only finitely many compatible upper bounds on the positive pair
$(\epsilon,\delta)$.

The continued vertical ovals are invariant under
$R_x(x,y)=(-x,y)$, and the horizontal ovals under
$R_y(x,y)=(x,-y)$, by Lemma~\ref{lem:continued-symmetry}.  Both
reflections leave $H_i$ invariant.  Functoriality of pullback gives
\[
 R_x^*\Phi^*\omega_{i-1}=R_y^*\Phi^*\omega_{i-1}
 =\Phi^*\omega_{i-1},\qquad
 R_x^*\alpha_U=\alpha_U,\quad R_y^*\alpha_V=\alpha_V.
\]
Consequently, on the entire vertical annuli the periods of
$\Phi^*\omega_{i-1}$ and $\alpha_U$ vanish exactly.  On the entire
horizontal annuli the periods of $\Phi^*\omega_{i-1}$ and $\alpha_V$
vanish exactly.  Padding changes the old copied periods by a small
amount, but does not change any of these fresh-source identities.

\paragraph{Auxiliary births and the global form.}
With the padded potentials now fixed, complete the auxiliary forms for
each remaining label $t>i$.  Start with the one-coordinate pullback of
$\Omega_{i-1}^{(t)}$, whose selected simple zeros have survived.  On a
new outer annulus of $z^2/2+U_i(w)$ or $z^2/2+V_i(w)$, this pulled form
has identically zero period: it is invariant under $w\mapsto-w$,
which reverses the oval orientation.  In the span of \eqref{eq:mechanical-birth-modes}, choose a fresh birth
form with one fewer selected simple period zero than the dimension of
that span. Add it to the pulled form with one sufficiently small nonzero
scalar. The sum retains every selected inherited zero and supplies the
chosen new simple zeros on an outer interval.
For either padded potential $W=U_i$ or $V_i$, choose the new outer
interval above all real critical values of $W$ and all values of
$z^2/2+W(w)$ on the previously selected compact annular neighborhoods.
Such an interval exists: these neighborhoods are finite in number and
compact, and $W$ has positive leading coefficient. Each mechanical level
there is a unique outer oval, so the new annular neighborhoods are disjoint
from those previously selected. The degree $2n_i$ and positive leading
coefficient give the required rank, however small the padding coefficient.
The forms for different future stages are chosen independently; these
choices change neither $H_i$ nor the current global form.

Finally define the single global form
\begin{equation}\label{eq:mechanical-global-form}
 \omega_i=\Phi^*\omega_{i-1}
                  +\lambda_V\alpha_V+\lambda_U\alpha_U.
\end{equation}
Choose two nonzero scalars sufficiently small to retain all four old
global copies.  On a fresh vertical annulus its period is exactly
$\lambda_V\oint\alpha_V$, and on a fresh horizontal annulus it is
exactly $\lambda_U\oint\alpha_U$.  Thus all new zeros coexist for this
same coefficient vector.  Each newly chosen scalar is fixed before
the next stage: the resulting positive, possibly very small, simple-zero
margins are added to the induction data.  No uniform margin over the
future degree or over all polynomial budgets is required.  This completes the finite induction.

\paragraph{Degrees and number of zeros.}
For an auxiliary monomial with $w$-coefficient degree $d$, the vertical
or horizontal source degree is at most
\[
 (2p+1)+2d+1=2(p+d+1).
\]
The final auxiliary budget for label $i$ is $p+d\leq n_i-2$, so each
source has degree at most $2n_i-2$.  If
$\deg\omega_{i-1}\leq2n_{i-1}-1$, then
\[
 \deg\Phi^*\omega_{i-1}
 \leq2(2n_{i-1}-1)+1
 =4n_{i-1}-1\leq2n_i-1.
\]
Also $\deg dH_i=2n_i-1$.  These estimates prove every required degree
bound for~\eqref{eq:mechanical-global-form} and the final field.

At stage $i$, put $s=n_{i-1}$, $N=n_i=2s+\eta_i$. For an auxiliary
birth with $k$ remaining folds let $A=2^k$,
$d=\lfloor s/A\rfloor$, and $e=s-Ad=e_{i,k}$.  Its actual birth
half-degree is $d$.  The allowed modes in
\eqref{eq:mechanical-birth-modes} number
\[
 M_k=\sum_{q=0}^{d-1}\{N-A(2q+1)\}
     =d(N-Ad).
\]
All summands are positive: at the largest $q$, their value is
$2e+\eta_i+A\geq1$.  In particular $M_k\geq1$.  The $A$ real
copies contribute $A(M_k-1)$ simple zeros.  The exact identity
\[
 A M_k=(s-e)(s+\eta_i+e)
       =s(s+\eta_i)-e^2-\eta_i e
\]
shows that four source copies give
$4\sum_{k=0}^{i-1}A(M_k-1)=K_i$, as in
\eqref{eq:mechanical-birth}.  This also proves $K_i\geq0$ and
integrality. The selected-zero counts obey
$L_0=0$, $L_i=4L_{i-1}+K_i$, yielding
\eqref{eq:mechanical-total}.

Write $\omega_r=P\,dx+Q\,dy$.  Lemma~\ref{lem:melnikov}, applied to
the finite disjoint configuration, supplies one sufficiently small
nonzero $\mu$ for which
\[
 X_\mu=( (H_r)_y-\mu Q,\,-(H_r)_x+\mu P )
\]
has a hyperbolic limit cycle near each selected simple period zero.
This is one real polynomial vector field, of degree at most $2n-1$.
\end{proof}

